\chapter{Options Pricing: Black-Scholes and Beyond}
\label{ch:options}

\begin{companioncode}{quant-options}
Code: \texttt{crates/quant-options/}\\
Dependencies: \crate{quant-core}, \crate{quant-stochastic}, \crate{thiserror}\\
Run: \texttt{cargo run -p quant-options --example greeks}
\end{companioncode}

\section{Learning Objectives}

This chapter builds the full Black-Scholes options toolkit on top of the
closed-form pricing primitives from \cref{ch:stochastic}: the Greeks
(Delta, Gamma, Vega, Theta, Rho) analytically and by finite difference,
and implied volatility via Newton's method with a bisection fallback.
After completing it you can:

\begin{itemize}
    \item State the Black-Scholes PDE and the risk-neutral derivation
    \item Compute all five first- and second-order Greeks in closed form
    \item Recover the same Greeks by finite difference for any pricing
      model (not just Black-Scholes)
    \item Solve for implied volatility from a market price using Newton's
      method, and explain why a bisection fallback is needed for deep
      in/out-of-the-money options
    \item Reproduce a synthetic volatility smile and explain its shape
\end{itemize}

\begin{keyinsight}
The Greeks are the partial derivatives of the option price with respect to
its inputs. Delta-hedging (holding $-\Delta$ shares against a long call)
eliminates first-order price risk; Gamma controls the residual hedging
error. Implied volatility inverts the Black-Scholes formula to recover
the market's view of future variance --- the ``volatility smile'' is the
single most studied phenomenon in derivatives.
\end{keyinsight}

\section{The Black-Scholes PDE}

\marginnote{%
\textbf{Black-Scholes PDE}\\[0.3em]
$\dfrac{\partial V}{\partial t} + \dfrac{1}{2}\sigma^2 S^2
\dfrac{\partial^2 V}{\partial S^2} + r S \dfrac{\partial V}{\partial S}
- r V = 0$
}

The Black-Scholes model assumes the underlying follows geometric Brownian
motion (see \cref{ch:stochastic}):
\[
dS_t = \mu S_t \, dt + \sigma S_t \, dW_t.
\]
By holding a continuously rebalanced portfolio of the option and
$-\Delta$ shares of stock, the residual risk is eliminated to first order.
No-arbitrage then requires the portfolio to earn the risk-free rate,
which yields the Black-Scholes PDE:
\[
\frac{\partial V}{\partial t} + \frac{1}{2}\sigma^2 S^2
\frac{\partial^2 V}{\partial S^2} + r S \frac{\partial V}{\partial S}
- r V = 0.
\]
Solving this PDE with the terminal payoff $\max(S_T - K, 0)$ for a call
(or $\max(K - S_T, 0)$ for a put) gives the closed forms
\[
C = S_0 \Phi(d_1) - K e^{-rT} \Phi(d_2), \qquad
P = K e^{-rT} \Phi(-d_2) - S_0 \Phi(-d_1),
\]
with
\[
d_1 = \frac{\ln(S_0/K) + (r + \frac{1}{2}\sigma^2)T}{\sigma\sqrt{T}},
\quad d_2 = d_1 - \sigma\sqrt{T}.
\]
These are the formulas we reuse from \crate{quant-stochastic::blackscholes}.

\section{The Greeks as Partial Derivatives}

\marginnote{%
\textbf{Greeks}\\[0.3em]
$\Delta = \dfrac{\partial C}{\partial S}$ \quad (spot)\\
$\Gamma = \dfrac{\partial^2 C}{\partial S^2}$ \quad (curvature)\\
$\mathcal{V} = \dfrac{\partial C}{\partial \sigma}$ \quad (vol)\\
$\Theta = \dfrac{\partial C}{\partial t}$ \quad (time)\\
$\rho  = \dfrac{\partial C}{\partial r}$ \quad (rate)
}

The Greeks are the sensitivities of the option price to its inputs. They
are the building blocks of hedging and risk management:

\begin{description}
    \item[Delta ($\Delta$)] $\partial C / \partial S$. The hedge ratio:
      holding $-\Delta$ shares against a long option removes first-order
      spot risk. A call has $\Delta \in (0, 1)$; a put $\Delta \in (-1,
      0)$.
    \item[Gamma ($\Gamma$)] $\partial^2 C / \partial S^2$. The rate of
      change of Delta. High Gamma means the hedge must rebalance
      frequently --- Gamma is largest at-the-money and shrinks in the
      wings. Identical for calls and puts.
    \item[Vega ($\mathcal{V}$)] $\partial C / \partial \sigma$. The
      sensitivity to volatility. Identical for calls and puts.
    \item[Theta ($\Theta$)] $\partial C / \partial t$. The time decay.
      Long options lose value as expiry approaches, so $\Theta < 0$ for
      long calls and puts.
    \item[Rho ($\rho$)] $\partial C / \partial r$. The sensitivity to the
      risk-free rate. Calls gain with $r$ (positive $\rho$); puts lose
      (negative $\rho$).
\end{description}

\section{Analytical Greeks}

\marginnote{%
\textbf{Analytical Greeks (call)}\\[0.3em]
$\Delta = \Phi(d_1)$\\
$\Gamma = \dfrac{\phi(d_1)}{S_0 \sigma \sqrt{T}}$\\
$\mathcal{V} = S_0 \phi(d_1) \sqrt{T}$\\
$\Theta = -\dfrac{S_0 \phi(d_1) \sigma}{2\sqrt{T}} - r K e^{-rT} \Phi(d_2)$\\
$\rho = K T e^{-rT} \Phi(d_2)$
}

Closed forms for the five Greeks follow from differentiating the
Black-Scholes price. Let $\phi(x) = (2\pi)^{-1/2} e^{-x^2/2}$ be the
standard normal PDF (the derivative of $\Phi$). We implement it inline
in \crate{quant-options::greeks}:

\begin{lstlisting}[language=Rust]
/// Standard normal PDF phi(x) = (1/sqrt(2 pi)) exp(-x^2 / 2).
pub fn normal_pdf(x: f64) -> f64 {
    const INV_SQRT_2PI: f64 = 0.398_942_280_401_432_7;
    INV_SQRT_2PI * (-0.5 * x * x).exp()
}
\end{lstlisting}

The analytical Greeks then drop out:

\begin{lstlisting}[language=Rust]
pub fn delta(s0, k, r, sigma, t, is_call: bool) -> f64 {
    let d1_val = d1(s0, k, r, sigma, t);
    let n_d1 = normal_cdf(d1_val);
    if is_call { n_d1 } else { n_d1 - 1.0 }
}

pub fn gamma(s0, k, r, sigma, t) -> f64 {
    let d1_val = d1(s0, k, r, sigma, t);
    normal_pdf(d1_val) / (s0 * sigma * t.sqrt())
}

pub fn vega(s0, k, r, sigma, t) -> f64 {
    let d1_val = d1(s0, k, r, sigma, t);
    s0 * normal_pdf(d1_val) * t.sqrt()
}
\end{lstlisting}

Theta and Rho differ between call and put; the call versions are
\[
\Theta_c = -\frac{S_0 \phi(d_1) \sigma}{2\sqrt{T}} - r K e^{-rT} \Phi(d_2),
\qquad
\rho_c = K T e^{-rT} \Phi(d_2).
\]
The put versions follow from put-call parity:
$\Theta_p = \Theta_c + r K e^{-rT}$ and
$\rho_p = -K T e^{-rT} \Phi(-d_2) = \rho_c - K T e^{-rT}$.

\section{Numerical Greeks by Finite Difference}

\marginnote{%
\textbf{Central difference}\\[0.3em]
$f'(x) \approx \dfrac{f(x+h) - f(x-h)}{2 h}$ \\[0.3em]
Error: $O(h^2)$
}

Analytical Greeks require a closed form. For models without a closed form
(local volatility, Heston, Monte Carlo) the same sensitivities are
recovered by finite difference. The central difference
\[
f'(x) \approx \frac{f(x+h) - f(x-h)}{2h}
\]
is second-order accurate: the symmetric form cancels the $O(h)$
truncation term, leaving $O(h^2)$. For Gamma (a second derivative), the
central second difference
\[
f''(x) \approx \frac{f(x+h) - 2 f(x) + f(x-h)}{h^2}
\]
is likewise $O(h^2)$.

\begin{lstlisting}[language=Rust]
pub fn delta_fd(s0, k, r, sigma, t, is_call, h) -> f64 {
    let up = price(s0 + h, k, r, sigma, t, is_call);
    let down = price(s0 - h, k, r, sigma, t, is_call);
    (up - down) / (2.0 * h)
}

pub fn gamma_fd(s0, k, r, sigma, t, h) -> f64 {
    let up = price(s0 + h, k, r, sigma, t, true);
    let mid = price(s0, k, r, sigma, t, true);
    let down = price(s0 - h, k, r, sigma, t, true);
    (up - 2.0 * mid + down) / (h * h)
}
\end{lstlisting}

\begin{keyinsight}
Finite difference is a sanity check on the analytical Greeks. If the two
disagree by more than a few basis points, one of them is wrong. With
$h = 10^{-4}$ for spot bumps and $h = 10^{-3}$ for the second derivative,
the analytical and numerical Greeks agree to $10^{-4}$ or better --- the
truncation and round-off errors are well balanced.
\end{keyinsight}

Theta uses a \emph{forward} difference, not central, because we cannot
step to $t - h < 0$ for a short-dated option:
\[
\Theta \approx \frac{C(t - h) - C(t)}{h}.
\]
The sign is the ``desk convention'': Theta is the negative of
$\partial C / \partial T$ (time to maturity), since long options lose
value as calendar time advances.

\section{Implied Volatility}

\marginnote{%
\textbf{Newton's method}\\[0.3em]
$\sigma_{n+1} = \sigma_n - \dfrac{C(\sigma_n) - C_{\text{mkt}}}{\mathcal{V}(\sigma_n)}$
}

Given a market price $C_{\text{mkt}}$, the implied volatility is the
$\sigma$ such that $C(S, K, r, \sigma, T) = C_{\text{mkt}}$. Because the
Black-Scholes price is monotone increasing in $\sigma$ (Vega is positive),
there is a unique solution. Newton's method is the textbook choice
because Vega is the exact derivative:
\[
\sigma_{n+1} = \sigma_n - \frac{C(\sigma_n) - C_{\text{mkt}}}{\mathcal{V}(\sigma_n)}.
\]
Newton converges quadratically when the initial guess is reasonable and
the option is not deep in/out of the money. When the option is deep ITM
or OTM, Vega collapses toward zero (the option price is intrinsic and
barely moves with $\sigma$), and the Newton step becomes ill-conditioned.
A bisection fallback on $[\sigma_{\min}, \sigma_{\max}]$ is then robust:
bisection is linear but guaranteed.

\begin{lstlisting}[language=Rust]
let newton_step = if v.abs() > VEGA_FLOOR {
    sigma - diff / v
} else {
    f64::NAN  // hand off to bisection
};

let next = if newton_step.is_finite()
    && newton_step > lo
    && newton_step < hi
{
    newton_step
} else {
    0.5 * (lo + hi)  // bisection
};
\end{lstlisting}

The initial guess uses the Brenner-Subrahmanyam (1988) approximation
$\sigma_0 \approx \sqrt{2\pi / T} \cdot C / S_0$, which is exact for
at-the-money options and bracketed by bisection if it diverges.

\section{Put-Call Parity and Implied Vol}

\marginnote{%
\textbf{Put-call parity}\\[0.3em]
$C - P = S_0 - K e^{-rT}$
}

Put-call parity says $C - P = S_0 - K e^{-rT}$. The implied vol that
reproduces a call price must therefore equal the implied vol that
reproduces the corresponding put price. The solver always inverts the
call formula; for a put we translate the put price into the equivalent
call price via $C = P + S_0 - K e^{-rT}$ and proceed. The IV is invariant
across calls and puts of the same strike and maturity.

\section{Results and Analysis}

\subsection{Greeks for an ATM Call}

Pricing an at-the-money call ($S_0 = K = 100$, $r = 0.05$, $\sigma = 0.2$,
$T = 1$) gives the call price $C = 10.4506$ (matching the Black-Scholes
benchmark from \cref{ch:stochastic}) and the Greeks:

\begin{center}
\begin{tabular}{lrrr}
\toprule
Greek & Analytical & Finite diff. & $|$diff$|$ \\
\midrule
$\Delta_c$ &  0.636831 &  0.636836 & $5.24 \times 10^{-6}$ \\
$\Delta_p$ & -0.363169 & -0.363164 & $5.24 \times 10^{-6}$ \\
$\Gamma$   &  0.018762 &  0.018763 & $1.34 \times 10^{-6}$ \\
$\mathcal{V}$ & 37.524035 & 37.523872 & $1.63 \times 10^{-4}$ \\
$\Theta_c$ & -6.414028 & -6.414142 & $1.15 \times 10^{-4}$ \\
$\Theta_p$ & -1.657881 & -1.657983 & $1.03 \times 10^{-4}$ \\
$\rho_c$   & 53.232483 & --- & --- \\
$\rho_p$   & -41.890459 & --- & --- \\
\bottomrule
\end{tabular}
\end{center}

The analytical and finite-difference Greeks agree to $10^{-4}$ or better
with $h = 10^{-4}$ for spot/vol bumps and $h = 10^{-3}$ for Gamma. Theta
matches to $10^{-4}$ with a forward difference in $t$.

\subsection{Implied Volatility Recovery}

Recovering $\sigma$ from the ATM call price $C = 10.4506$:

\begin{lstlisting}[language=bash]
$ cargo run -p quant-options --example implied_vol
Recovery (ATM call):
  Market price  = 10.450575
  Implied vol   = 0.20000000
  True vol      = 0.20000000
  |iv - true|   = 3.54e-14
\end{lstlisting}

The implied vol recovers the true $\sigma = 0.2$ to $\sim 10^{-14}$ ---
far below the typical bid-ask spread. The same IV is recovered from the
put price, as put-call parity demands.

\subsection{Volatility Smile}

A synthetic smile is generated by pricing calls under
$\sigma(K) = \sigma_0 + a \cdot (\ln(K/S_0))^2$ with $a = 0.5$, so the
ATM vol is $\sigma_0$ and the wings have higher vol (the classic equity
smile). Inverting the BS formula at each strike recovers the smile:

\begin{center}
\begin{tabular}{rrrr}
\toprule
$K$ & moneyness $K/S_0$ & market price & IV \\
\midrule
 70 & 0.700 & 33.981835 & 0.263609 \\
 80 & 0.800 & 24.965341 & 0.224897 \\
 90 & 0.900 & 16.851425 & 0.205550 \\
 95 & 0.950 & 13.390174 & 0.201316 \\
100 & 1.000 & 10.450575 & 0.200000 \\
105 & 1.050 &  8.068580 & 0.201190 \\
110 & 1.100 &  6.219914 & 0.204542 \\
120 & 1.200 &  3.823180 & 0.216621 \\
130 & 1.300 &  2.578807 & 0.234418 \\
\bottomrule
\end{tabular}
\end{center}

The IV is symmetric around the ATM strike in log-moneyness, exactly as
the model prescribes. In real markets the smile is not symmetric: equity
smiles are skewed (puts bid up by crash protection), FX smiles are more
symmetric, and rates smiles can be monotone. Fitting these is the job of
local-volatility, stochastic-volatility, and SVI parameterisations.

\begin{keyinsight}
The smile is the market's fingerprint. If the Black-Scholes model were
literally true, IV would be a flat line at $\sigma$. Real markets quote
options at many strikes with different IVs, and the shape of that curve
encodes the market's view of tail risk. Modelling the smile is the bridge
between the Black-Scholes world and the post-1987 derivatives industry.
\end{keyinsight}

\section{Exercises}

\begin{enumerate}
    \item \textbf{Black-Scholes with dividends.} Extend the BS formula to
      a continuous dividend yield $q$: replace $S_0$ with $S_0 e^{-qT}$
      everywhere, and re-derive the Greeks. Verify Delta becomes
      $e^{-qT} \Phi(d_1)$ for the call.

    \item \textbf{Local volatility.} Implement a simple local-volatility
      pricer where $\sigma(S, t)$ is a deterministic function of spot and
      time (e.g. $\sigma(S, t) = \sigma_0 + \alpha (S - S_0)^2$). Price an
      option by explicit Euler on a spot-time grid and compare the smile
      to the IV surface.

    \item \textbf{Heston model.} Implement the Heston stochastic-volatility
      SDE $dv_t = \kappa(\theta - v_t) dt + \xi \sqrt{v_t} dW^{(v)}_t$
      with Monte Carlo. Generate call prices for a strip of strikes and
      invert to recover the smile. Compare to the equity smile shape.

    \item \textbf{Vega-clipped bisection.} For an extreme deep-ITM call
      (e.g. $S_0 = 1000$, $K = 10$) show that the BS price equals the
      intrinsic value to machine precision, so IV is genuinely
      unrecoverable. Modify the solver to return \texttt{Ok(sigma\_min)}
      with a warning flag rather than attempting to recover.

    \item \textbf{Smile fitting.} Download a real option chain (e.g.
      SPX) and fit a quadratic-in-log-moneyness smile
      $\sigma(K) = \beta_0 + \beta_1 \ln(K/S_0) + \beta_2 (\ln(K/S_0))^2$
      via least squares. Report the residual and the skew coefficient
      $\beta_1$.

    \item \textbf{Newton convergence rate.} Instrument \func{implied\_vol}
      to print $|C(\sigma_n) - C_{\text{mkt}}|$ at each iteration. Verify
      quadratic convergence for ATM options (the error roughly squares
      each step) and linear convergence once bisection takes over.
\end{enumerate}

\section{Key Takeaways}

\begin{itemize}
    \item \textbf{Black-Scholes PDE} prices any payoff via continuous
      delta-hedging; the closed-form call and put follow from solving
      the PDE with the terminal payoff.
    \item \textbf{Greeks} are the partial derivatives of the price.
      Delta is the hedge ratio; Gamma the curvature; Vega the vol
      sensitivity; Theta the time decay; Rho the rate sensitivity. Gamma
      and Vega are identical for calls and puts.
    \item \textbf{Finite difference} recovers the Greeks for any pricing
      model, not just Black-Scholes. Central differences are $O(h^2)$;
      forward differences are needed for Theta (cannot step to $t < 0$).
    \item \textbf{Implied volatility} inverts the BS formula via Newton's
      method, with a bisection fallback when Vega collapses (deep ITM/OTM).
    \item \textbf{Put-call parity} implies the IV from a call equals the
      IV from the corresponding put --- a basic arbitrage check that
      real-market quotes must satisfy.
    \item \textbf{The volatility smile} is the market-implied IV curve
      across strikes; modelling it is the bridge from Black-Scholes to
      post-1987 derivatives practice.
\end{itemize}

\vspace{1em}
\noindent\textbf{Next chapter}: Portfolio optimization --- Markowitz
mean-variance frontier, the efficient frontier, and the capital asset
pricing model.